%------------------------------------------------------------------------------%
\begin{exercises}

% 1 e 2 Stewart pt pg 633
%------------------------------------------------------------------------------%
\begin{exercise}
  Responda as questões.
  \begin{alphalist}
    \item O que é uma sequência?
    \item O que significa dizer que \(\lim_{n\to\infty} a_n = 5 \)?
    \item O que significa dizer que \(\lim_{n\to\infty} a_n = 0 \)?
    \item O que é uma sequência convergente? Dê dois exemplos.
    \item O que é uma sequência divergente? Dê dois exemplos.
  \end{alphalist}
\end{exercise}

% 01 - Thomas pg 10
%------------------------------------------------------------------------------%
\begin{exercise}
Dada a fórmula para o $n$-ésimo termo, $a_n$, de uma sequência $(a_n)$,
encontre os valores de $a_1$, $a_2$, $a_3$ e $a_4$.
\[
  a_n = \frac{1-n}{n^2}
\]
\begin{answer}
\begin{align*}
  a_1 & = \dfrac{\;1-n\;}{n^2}\evalat{n=1}{} = \dfrac{\;1-1\;}{1^2} = 0              \\[3mm]
  a_2 & = \dfrac{\;1-n\;}{n^2}\evalat{n=2}{} = \dfrac{\;1-2\;}{2^2} = -\dfrac{1}{4}  \\[3mm]
  a_3 & = \dfrac{\;1-n\;}{n^2}\evalat{n=3}{} = \dfrac{\;1-3\;}{3^2} = -\dfrac{2}{9}  \\[3mm]
  a_4 & = \dfrac{\;1-n\;}{n^2}\evalat{n=4}{} = \dfrac{\;1-4\;}{4^2} = -\dfrac{3}{16}
\end{align*}
\end{answer}
\end{exercise}

% 02-06 - Thomas pg 10
% 19-22 - Stewart pt pg 633
%------------------------------------------------------------------------------%
\begin{exercise}
  Dada a fórmula para o $n$-ésimo termo $a_n$,
  encontre os valores dos quatro primeiros termos da sequência $(a_n)$.
  \begin{mathlist}[2]
    \item a_n = \frac{1}{n!}
    \item a_n = \frac{(-1)^{n+1}}{2n-1}
    \item a_n = 2 + (-1)^n
    \item a_n = \frac{2^n}{2^{n+1}}
    \item a_n = \frac{2^n-1}{2^n}
    %
    \item a_n = \frac{3n}{1+6n}
    \item a_n = 2 + \frac{(-1)^n}{n}
    \item a_n = 1 + \left(-\frac{1}{2}\right)^n
    \item a_n = 1 + \frac{10^n}{9^n}
  \end{mathlist}
\end{exercise}

% 10 - Thomas pg 10
%--------------------------------------------------------------------------------------------------%
\begin{exercise}
Dada um ou dois termos iniciais de uma sequência, bem como uma fórmula de recursão para os termos
remanescentes. Escreva os dez termos iniciais da sequência.
\[
  a_1 = -2\qquad a_{n+1} = \frac{n a_n}{\;n+1\;}
\]
\begin{answer}
\begin{align*}
  a_1    & = -2                                                                                 \\[3mm]
  a_2    & = \frac{1 a_1}{\;1+1\;} = \frac{    -2}           { 2} = -1                          \\[3mm]
  a_3    & = \frac{2 a_2}{\;2+1\;} = \frac{2 (-1)}           { 3} = -\frac{2}{3}                \\[3mm]
  a_4    & = \frac{3 a_3}{\;3+1\;} = \frac{3 (-\sfrac{2}{3})}{ 4} = -\frac{2}{4} = -\frac{1}{2} \\[3mm]
  a_5    & = \frac{4 a_4}{\;4+1\;} = \frac{4 (-\sfrac{1}{2})}{ 5} = -\frac{2}{5}                \\[3mm]
  a_6    & = \frac{5 a_5}{\;5+1\;} = \frac{5 (-\sfrac{2}{5})}{ 6} = -\frac{2}{6} = -\frac{1}{3} \\[3mm]
  a_7    & = \frac{6 a_6}{\;6+1\;} = \frac{6 (-\sfrac{1}{3})}{ 7} = -\frac{2}{7}                \\[3mm]
  a_8    & = \frac{7 a_7}{\;7+1\;} = \frac{7 (-\sfrac{2}{7})}{ 8} = -\frac{2}{8} = -\frac{1}{4} \\[3mm]
  a_9    & = \frac{8 a_8}{\;8+1\;} = \frac{8 (-\sfrac{1}{4})}{ 9} = -\frac{1}{9}                \\[3mm]
  a_{10} & = \frac{9 a_9}{\;9+1\;} = \frac{9 (-\sfrac{1}{9})}{10} = -\frac{1}{10}               \\[3mm]
\end{align*}
\end{answer}
\end{exercise}

% 07-09 e 11-12 - Thomas pg 10
%------------------------------------------------------------------------------%
\begin{exercise}
  Dada a relação de recorrência e os valores iniciais necessários da sequência $(a_n)$,
  calcule seus cinco primeiros termos.
  \begin{mathlist}
  \item a_1 = 1             \) \tabto{40mm} \( a_{n+1} = a_n + \left(\frac{1}{2}\right)^2 \)
  \item a_1 = 1             \) \tabto{40mm} \( a_{n+1} = \frac{a_n}{n+1} \)
  \item a_1 = 2             \) \tabto{40mm} \( a_{n+1} = (-1)^{n+1}\frac{a_n}{2} \)
  \item a_1 = a_2 = 1       \) \tabto{40mm} \( a_{n+2} = a_{n+1} + a_n  \)
  \item a_1 = 2,\; a_2 = -1 \) \tabto{40mm} \( a_{n+2} = \frac{a_{n+1}}{a_n} \)
  \end{mathlist}
\end{exercise}

% 19 - Thomas pg 10
%--------------------------------------------------------------------------------------------------%
\begin{exercise}
Encontre uma fórmula para o $n$-ésimo termo da sequência
dos quadrados dos inteiros positivos menos $1$
\[ \left(0,\, 3,\, 8,\, 15,\, 24,\, \dots \right) \]
\begin{answer}
\[
  a_n = n^2-1
\]
Verificando
\begin{align*}
  a_1 & = (n^2-1)\evalat{n=1}{} = 1^2-1 = 1^2-1 = 0  \\[3mm]
  a_2 & = (n^2-1)\evalat{n=2}{} = 2^2-1 = 2^2-1 = 3  \\[3mm]
  a_3 & = (n^2-1)\evalat{n=3}{} = 3^2-1 = 3^2-1 = 8  \\[3mm]
  a_4 & = (n^2-1)\evalat{n=4}{} = 4^2-1 = 4^2-1 = 15 \\[3mm]
  a_5 & = (n^2-1)\evalat{n=5}{} = 5^2-1 = 5^2-1 = 24
\end{align*}
\emph{Obs:}
sem a informação de que a sequência é formada por ``Quadrados dos inteiros positivos menos~$1$'',
não existe uma solução única para o problema.
A expressão polinomial a seguir é outra solução possível para gerar os valores listados
\[
  a_n = -\frac{1}{3}n^4 +\frac{10}{3}n^3 -\frac{32}{3}n^2 +\frac{50}{3}n -9
\]
\end{answer}
\end{exercise}

% 13, 15, 17, 20, 21, 26 - Thomas pg 10
% 13-18 Stewart pt pg 633
%------------------------------------------------------------------------------%
\begin{exercise}
  Encontre uma fórmula explicita para o $n$-ésimo termo de cada sequência.
  \begin{alphalist}
  \item Números 1 alternando o sinal
    \[ \left( 1, -1, 1, -1,      \dots \right) \]
  \item Quadrados dos inteiros positivos alternando os sinais
    \[ \left( 1, -4, 9, -16, 25, \dots \right) \]
  \item Potências de $2$ divididas por múltiplos de $3$
    \[
      \left( \frac{1}{9},
             \frac{2}{12},
             \frac{2^2}{15},
             \frac{2^3}{18},
             \dots
      \right)
    \]
  \item Inteiros começando com $-3$
    \[ \left( -3, -2, -1, 0, 1, \dots \right) \]
  \item Um inteiro positivo impar sim um não
    \[ \left( 1, 5, 9, 13, 17,  \dots \right) \]
  \item
    \( \left( 1,
      \frac{1}{3},
      \frac{1}{5},
      \frac{1}{7},
      \frac{1}{9},
      \dots \right)
    \)
  \item
    \( \left( 1,
     -\frac{1}{3},
      \frac{1}{9},
     -\frac{1}{27},
      \frac{1}{81},
      \dots \right)
    \)
  \item
    \( \left( -3, 2,
     -\frac{4}{3},
      \frac{8}{9},
     -\frac{16}{27},
      \dots \right)
    \)
  \item \( \left( 5, 8, 11, 14, 17, \dots \right)\)
  \item
    \( \left(
      \frac{1}{2},
     -\frac{4}{3},
      \frac{9}{4},
     -\frac{16}{5},
      \frac{25}{6},
      \dots \right)
    \)
  \item \( \left( 1, 0, -1, 0, 1, 0, -1, 0, \dots \right)\)
  \end{alphalist}
\end{exercise}

% Ex 35 pg 10
%--------------------------------------------------------------------------------------------------%
\begin{exercise}
A sequência $(a_n)$ converge ou diverge? Encontre seu limite se for convergente.
\[
  a_n = 1 + (-1)^n
\]
\begin{answer}
Os termos dessa sequência assumem os valores 0 e 2 alternadamente, portanto ela não converge.
\end{answer}
\end{exercise}

% Ex 36 pg 10
%--------------------------------------------------------------------------------------------------%
\begin{exercise}
A sequência $(a_n)$ converge ou diverge? Encontre seu limite se for convergente.
\[
  a_n = (-1)^n \left( 1 - \frac{1}{n} \right)
\]
\begin{answer}
Iniciamos calculando o limite
\[
  \lim_{n\to\infty} \left( 1 - \frac{1}{n} \right) = \lim  1 - \lim \frac{1}{n} = 1 - 0 = 1
\]
Tomando os termos pares e impares separadamente temos
\begin{align*}
	\lim_{n\to\infty} a_{2n} & = \lim_{n\to\infty} \left[(-1)^{2n}   \left( 1 - \frac{1}{2n}   \right)\right] \\
	                         & = \lim_{n\to\infty} \left( 1 - \frac{1}{2n} \right)                            \\
	                         & = \lim_{k\to\infty} \left( 1 - \frac{1}{k} \right) = 1
\end{align*}
\begin{align*}
	\lim_{n\to\infty} a_{2n+1} & = \lim_{n\to\infty} \left[(-1)^{2n+1}   \left( 1 - \frac{1}{2n+1}   \right)\right] \\
	                           & = \lim_{n\to\infty} -\left( 1 - \frac{1}{2n} \right)                               \\
	                           & = \lim_{k\to\infty} -\left( 1 - \frac{1}{k} \right) = -1
\end{align*}
Como cada subsequência converge para um valor diferente e o limite é único quando existe concluímos
que a sequência diverge.
\end{answer}
\end{exercise}

% Seção 10.1 - Exercício 101 página 11
%------------------------------------------------------------------------------%
\begin{exercise}
Método de Newton: As seguintes sequências vêm da fórmula recursiva para o método de Newton,
\[
  x_{n+1} = x_n - \frac{f(x_n)}{\;f'(x_n)\;}
\]
As sequências convergem?
Em caso afirmativo, para qual valor?
Em cada caso, comece identificando a função $f$ que gera a sequência.
\begin{mathlist}
  \item x_0 = 1,\quad x_{n+1} = x_n - \dfrac{x_n^2-2}{2x_n} = \dfrac{x_n}{2}+\dfrac{1}{x_n}
  \item x_0 = 1,\quad x_{n+1} = x_n - \dfrac{\tan(x_n) - 1}{\sec^2(x_n)}
  \item x_0 = 1,\quad x_{n+1} = x_n - 1
\end{mathlist}
\begin{answer}
O método de Newton é usado para encontrar os zeros de funções reais, $f(x)=0$.
Se o valor inicial $x_0$ estiver próximo o suficiente da raiz temos a garantia
de convergência do método.
\begin{alphalist}
  \item A função usada nesse caso é $f(x)=x^2-2$ que possui raízes em $x=\pm\sqrt{2}$,
  portanto se a sequência convergir deve ser para um desses valores.
  Calculando os primeiros termos obtemos
  \begin{ceqn}
  \[
  \begin{array}{cc} \hline
  	   n     &       x_n       \\ \hline
  	   0     & 1,000\,000\,000 \\
  	   1     & 1,500\,000\,000 \\
  	   2     & 1,416\,666\,667 \\
  	   3     & 1,414\,215\,686 \\
  	   4     & 1,414\,213\,562 \\ \hline
  	\sqrt{2} & 1,414\,213\,562 \\ \hline
  \end{array}
  \]
  \end{ceqn}
  Observamos que a sequência converge para $1,414\,213\,562 \approx \sqrt{2}$.

  \item A função usada nesse caso é $f(x)=\tan(x)-1$ que possui raízes em
  $x=\sfrac{\pi}{4}+2\pi k$, para qualquer $k\in\Z$, portanto se a sequência
  convergir deve ser para um desses valores.
  \begin{ceqn}
  \[
  \begin{array}{cc} \hline
  	      n        &       x_n       \\ \hline
  	      0        & 1,000\,000\,000 \\
  	      1        & 0,837\,277\,868 \\
  	      2        & 0,788\,180\,293 \\
  	      3        & 0,785\,405\,918 \\
  	      4        & 0,785\,398\,163 \\ \hline
  	\sfrac{\pi}{4} & 0,785\,398\,163 \\ \hline
  \end{array}
  \]
  \end{ceqn}
  A sequência converge para $0,785\,3981\,635 \approx \sfrac{\pi}{4}$

  \item A função usada nesse caso é $f(x)=e^x$ que não possui raízes,
  portanto a série não pode convergir.
  A sequência diverge, assumindo os valores $1,0,-1,-2,-3,-4,-5,\dots$
\end{alphalist}
\end{answer}
\end{exercise}

% Exercícios 27-38 da página 10 do Thomas Vol. 2
%------------------------------------------------------------------------------%
\begin{exercise}
  Verifique se a sequência é convergente, caso seja calcule seu limite.
  \begin{mathlist}[2][series=ex-02-02-A]
     \item a_n = 2 + (0,1)^n
     \item a_n = \frac{n+(-1)^n}{n}
     \item a_n = \frac{1-2n}{1+2n}
     \item a_n = \frac{2n+1}{1-3\sqrt{n}}
     \item a_n = \frac{1-5n^4}{n^4+8n^3}
     \item a_n = \frac{n+3}{n^2+5n+6}
     \item a_n = \frac{n^2-2n+1}{n-1}
     \item a_n = \frac{1-n^3}{70-4n^2}
     \item a_n = 1+(-1)^n
   \end{mathlist}
   \begin{mathlist}[1][resume=ex-02-02-A]
     \item a_n = (-1)^n\left( 1-\frac{1}{n}\right)
     \item a_n = \left(\frac{n+1}{2n}\right) \left(1-\frac{1}{n}\right)
     \item a_n = \left(2-\frac{1}{2^n}\right) \left(3+\frac{1}{2^n}\right)
  \end{mathlist}
\begin{answer}
  \begin{alphalist}[2]
    \item $2$
    \item $1$
    \item $-1$
    \item Diverge para menos infinito
    \item $-5$
    \item $0$
    \item Diverge para mais infinito
    \item Diverge para mais infinito
    \item Diverge
    \item Diverge
    \item $\sfrac{1}{2}$
    \item $6$
  \end{alphalist}
\end{answer}
\end{exercise}

% 39-42 - Thomas pg 10
% 23-56 - Stewart pt pg 633
%------------------------------------------------------------------------------%
\begin{exercise}
  Verifique se a sequência converge ou não e calcule o limite das sequências convergentes.
  \begin{mathlist}[2]
  % 39 - 42
  \item a_n = \frac{(-1)^{n+1}}{2n-1}
  \item a_n = \left(-\frac{1}{2}\right)^n
  \item a_n = \frac{1}{(0,9)^n}
  % 23 - 30
  \item a_n = 1 - (0,2)^n
  \item a_n = \frac{n^3}{n^3+1}
  \item a_n = \frac{3+5n^2}{n+n^2}
  \item a_n = \frac{n^3}{n+1}
  \item a_n = \frac{3^{n+2}}{5^n}
  % 31 - 40
  \item a_n = \frac{(2n-1)!}{(2n+1)!}
  \end{mathlist}
\end{exercise}

% Exercícios 91-96 da página 11 do Thomas Vol. 2
%------------------------------------------------------------------------------%
\begin{exercise}
  Assuma que a sequência converge e encontre seu limite.
  \begin{mathlist}
    \item a_1 =  2 \) \tabto{30mm} \( a_{n+1} = \frac{72}{1+a_n}    \) % a
    \item a_1 = -1 \) \tabto{30mm} \( a_{n+1} = \frac{a_n+6}{a_n+2} \) % b
    \item a_1 = -4 \) \tabto{30mm} \( a_{n+1} = \sqrt{8+2a_n}       \) % c
    \item a_1 =  0 \) \tabto{30mm} \( a_{n+1} = \sqrt{8+2a_n}       \) % d
    \item a_1 =  5 \) \tabto{30mm} \( a_{n+1} = \sqrt{5a_n}         \) % e
    \item a_1 =  3 \) \tabto{30mm} \( a_{n+1} = 12 - \sqrt{a_n}     \) % f
    \item 2,\hfill
	    2+\frac{1}{2},\hfill
	    2+\frac{1}{2+\sfrac{1}{2}},\hfill
	    2+\frac{1}{2+\frac{1}{2+\sfrac{1}{2}}},\hfill
	    \dots \hfill % g
    \item \sqrt{1},\hfill
	    \sqrt{1+\sqrt{1}},\hfill
	    \sqrt{1+\sqrt{1+\sqrt{1}}},\hfill
	    \dots\hfill % h
  \end{mathlist}
  \begin{answer}
    \begin{alphalist}[3]
      \item $8$ % a
      \item $2$ % b
      \item $4$ % c
      \item $4$ % d
      \item $5$ % e
      \item $9$ % f
      \item $\sqrt{2}$ % g
      \item $\sfrac{1+\sqrt{5}}{2}$ % h
    \end{alphalist}
  \end{answer}
\end{exercise}

%------------------------------------------------------------------------------%
%\begin{exercise}
%\begin{answer}
%\end{answer}
%\end{exercise}

%------------------------------------------------------------------------------%
%\begin{exercise}
%\begin{mathlist}[2]
%  \item
%  \item
%  \item
%\end{mathlist}
%\begin{answer}
%  \begin{alphalist}[3]
%    \item
%    \item
%    \item
%  \end{alphalist}
%\end{answer}
%\end{exercise}

\end{exercises}
%------------------------------------------------------------------------------%
